\subsection{Auto Encoders}
Working with the full data representation is not always required and can be computationally expensive. In many domains, such as in image analysis or audio processing, the data often includes a large amount of features - pixels in images, or frequencies in audio signals. In these cases, it is possible and often beneficial to try and find a lower dimensional representation of the data. Encoders and decoders are very common in order to transmit and share for exaple images, video and audio - formats such as JPEG, H.264, and MP3 good examples, reducing the size of images and audio allowing for easier transportation.

Techniques such as Principal Component Analysis (PCA) can achieve the dimensionality reduction, by computing and projection the data into a lower dimensional hyperspace. However, PCA is a linear method and in some cases, the data is not linearly separable.

In contrast to PCA, autoencoders are a type of model that has two components to it; an \emph{encoder} and a \emph{decoder}. In contrast to PCA, the encoder and decoder are not neccesarily linear transformations, but can be any arbitrary function. The encoders maps the input from the original dimensional space to a lower dimensional space, usually denoted the \emph{latens space}.

$$
X_{latent} \in \mathbb{R}^H  = f_{encoder}(X_{input} \in \mathbb{R}^D)
$$

Where $X_{latent}$ is the latent representation of the input data $X_{input}$, $f_{encoder}$ is the encoder function and $H$ is the latent space dimension. The decoder then maps the latent representation back to the original dimensional space as $X_{reconstructed}$.

$$
X_{reconstructed} \in \mathbb{R}^D = f_{decoder}(X_{latent} \in \mathbb{R}^H)
$$

An ideal autoencoder would be able to fully reconstruct the original data from the latent representation received from the encoder, i.e., $X_{reconstructed} = f_{decoder}(f_{encoder}(X_{input}))$. However, most autoencoders will be lossy, and will not be able to fully reconstruct the original data.